\chapter{Diskussion}
\label{ch:diskussion}
% Reflektierender Teil -- oft unterschaetzt, aber genau das, was einen
% guten von einem rein beschreibenden Bericht unterscheidet.

\section{Herausforderungen}
% Die Speicher-Eskalation (Kap. 6) als zentrales Beispiel: naheliegende
% Erstdiagnosen trafen jeweils einen echten, aber nicht dominanten Faktor --
% erst die Frage "wie viele Objekte leben gleichzeitig" fuehrte zur Loesung.
% Das ist eine uebertragbare Erkenntnis, nicht nur eine Randnotiz.
Eine zentrale Herausforderung bei der Verarbeitung des vollständigen
Wikipedia-Korpus war der stark ansteigende Speicherbedarf. Während die
Verarbeitung kleinerer Datenmengen zunächst unauffällig verlief, zeigte sich
beim Skalieren auf den vollständigen Korpus, dass zuvor verwendete
In-Memory-Ansätze nicht ausreichend skalieren. Die Fehlersuche konzentrierte
sich zunächst auf einzelne speicherintensive Datenstrukturen und
Verarbeitungsschritte. Diese Faktoren trugen zwar zum Speicherverbrauch bei,
erklärten dessen Größenordnung jedoch nicht vollständig.

Entscheidend war letztlich nicht allein die Größe einzelner Objekte, sondern
die Anzahl der Objekte, die während der Verarbeitung gleichzeitig im Speicher
gehalten wurden. Dadurch konnte sich ein für einen einzelnen Artikel
unkritischer Speicherbedarf über mehrere Millionen Artikel zu einem erheblichen
Gesamtspeicherbedarf aufsummieren. Diese Beobachtung führte zu einer
grundlegenderen Betrachtung der Verarbeitungspipeline: Statt lediglich einzelne
Datenstrukturen zu optimieren, musste die Lebensdauer der Daten im Speicher
reduziert werden.

Daraus ergibt sich als übertragbare Erkenntnis, dass bei der Verarbeitung
großer Korpora nicht nur die Größe einzelner Datenstrukturen betrachtet werden
sollte. Ebenso relevant ist, wie viele dieser Strukturen gleichzeitig
existieren und zu welchem Zeitpunkt sie wieder freigegeben werden können. Für
WikiPod war diese Erkenntnis insbesondere für die Skalierung vom kleinen
Evaluationskorpus auf den vollständigen Wikipedia-Bestand von Bedeutung und
motivierte eine stärker streamingorientierte Verarbeitung, bei der
Zwischenergebnisse möglichst früh weiterverarbeitet beziehungsweise
persistiert werden, anstatt sie über den gesamten Verarbeitungslauf im
Arbeitsspeicher zu halten.

\section{Designentscheidungen im Rueckblick}
% Was wuerdet ihr nochmal genauso machen, was anders (z. B. frueher auf
% Streaming-Architektur setzen statt erst bei vollem Korpus zu merken,
% dass In-Memory-Ansaetze nicht skalieren).

Ein konkreter Punkt, den wir im Rückblick anders angehen würden, betrifft die Gewichtung des Scoring-Modells
aus Abschnitt~\ref{ch:dokumentenauswahl}. Die Messung auf dem vollen en.wikipedia-Korpus (Tabelle~\ref{tab:scoring-weights})
zeigt, dass \texttt{importance} trotz eines mit 0,150 vergleichsweise niedrigen Gewichts 82,61\,\% des
gewichteten Gesamtscores ausmacht, während \texttt{incoming\_links} trotz des höchsten Gewichts (0,250) nur
0,31\,\% beiträgt. Ursächlich ist, dass die logarithmische Dämpfung bei \texttt{importance} pro ausgehendem
Link und nicht auf den aggregierten Artikelwert angewendet wird -- der Rohwert wächst dadurch tendenziell mit
der Anzahl der ausgehenden Links und bleibt strukturell größer als bei den übrigen, einfach gedämpften
Signalen. Die konfigurierten Gewichte spiegeln damit nicht den tatsächlichen Einfluss der Signale auf die
Auswahl wider, obwohl genau das mit \texttt{config/prod.yaml} beziehungsweise \texttt{config/server.yaml}
beabsichtigt war.

Für eine zukünftige Iteration wären zwei Ansätze naheliegend: entweder \texttt{importance} durch die Anzahl
der berücksichtigten Links zu normalisieren (Mittelwert statt Summe der gedämpften Linkfrequenzen), oder alle
fünf Rohsignale vor der Gewichtung zusätzlich zu standardisieren (z.\,B. per Z-Score relativ zu ihrer eigenen
Verteilung im Korpus), sodass ein konfiguriertes Gewicht tatsächlich einen vergleichbaren Anteil am Endscore
bewirkt -- unabhängig von der ursprünglichen Skala des jeweiligen Signals. Beide Änderungen hätten jedoch
Rückwirkungen auf die bereits durchgeführten Selektionsläufe und wurden im Rahmen dieses Projekts aus
Zeitgründen nicht mehr umgesetzt.

\section{Grenzen der Loesung}
% Was deckt das System bewusst nicht ab (z. B. nur englische Wikipedia,
% keine Multi-Turn-Konversation, Kategorien-Signal nicht genutzt, ...).

Die entwickelte Lösung unterliegt mehreren Einschränkungen, die bei der
Einordnung der Ergebnisse berücksichtigt werden müssen. Eine grundlegende
Einschränkung ergibt sich aus der verwendeten Wissensbasis. WikiPod verarbeitet
ausschließlich Inhalte der englischsprachigen Wikipedia. Andere Sprachversionen
sowie zusätzliche externe Wissensquellen werden nicht berücksichtigt. Die
verfügbaren Informationen sind damit unmittelbar durch den Inhalt und den
Stand der lokal verwendeten Wikipedia-Kopie begrenzt.

Eine weitere Einschränkung betrifft die Auswahl der lokal vorgehaltenen
Dokumente. Aufgrund des vorgegebenen Speicherbudgets kann nur ein Teil des
verfügbaren Wikipedia-Bestands berücksichtigt werden. Die Auswahl basiert auf
dem entwickelten Scoring-Verfahren und stellt daher keine Garantie dafür dar,
dass für jede mögliche Nutzeranfrage der relevanteste Artikel im lokalen
Bestand enthalten ist. Wie die Analyse des Scoring-Modells gezeigt hat,
entspricht der tatsächliche Einfluss der verwendeten Signale zudem nicht in
allen Fällen ihrer konfigurierten Gewichtung. Insbesondere die Dominanz von
\texttt{importance} kann die Zusammensetzung der ausgewählten Teilmenge
beeinflussen.

Auch die Evaluation ist in ihrer Aussagekraft begrenzt. Der reguläre
Evaluationsdatensatz umfasst 20 Anfragen und wurde ebenso wie der zusätzliche
Edge-Case-Datensatz manuell zusammengestellt. Die gemessenen Retrieval-Metriken
beschreiben daher zuverlässig das Verhalten auf den untersuchten Anfragen,
lassen sich jedoch nicht ohne Weiteres auf die Gesamtheit möglicher
Nutzeranfragen übertragen. Insbesondere seltene Themen, weitere mehrdeutige
Bezeichnungen und Anfragen außerhalb des lokal verfügbaren Wissensbestands
können in den verwendeten Datensätzen nur eingeschränkt abgebildet werden.

Darüber hinaus wurde die Qualität der generierten Antworten nicht mit einer
stabilen quantitativen Metrik bewertet. Der im Rahmen der Evaluation erprobte
LLM-Judge zeigte bei einzelnen Anfragen inkonsistente Bewertungen und wurde
deshalb nicht als alleinige Grundlage für einen quantitativen Vergleich der
Antwortqualität verwendet. Der Vergleich der lokalen Sprachmodelle erlaubt
somit vor allem Aussagen über deren Generierungslatenz auf der Zielhardware,
nicht jedoch darüber, welches Modell insgesamt die höchste inhaltliche
Antwortqualität liefert.

Schließlich setzt die lokale Ausführung auf dem Raspberry Pi der Lösung
praktische Leistungsgrenzen. Während das semantische Retrieval im untersuchten
Aufbau mit vergleichsweise geringer Latenz durchgeführt werden konnte, benötigt
die lokale Antwortgenerierung deutlich mehr Zeit und stellt damit den
dominierenden Anteil der Antwortzeit dar. Die vollständig lokale Ausführung
ermöglicht zwar den Betrieb ohne bestehende Internetverbindung, geht auf der
verwendeten Hardware jedoch mit einem deutlichen Kompromiss zwischen
Modellgröße, Ressourcenbedarf und Generierungsgeschwindigkeit einher.
